% !TEX root = ../monograph.tex
\chapter[Common trajectories and stability]{A common trajectory and fixed-scale stability}
\label{ch10:trajectory}

The step must be used with a single schedule (Section~\ref{ch10:schedule})
and a consistent choice of activity arrays. There are two
separate problems. First one must tune each finite Wilson theory to the
same coarse inverse coupling. Then one must compare theories with different
ultraviolet cutoffs. Continuity of a finite step solves the first
problem, but does not by itself solve the second.

We first construct the numerical schedule and the tuned finite
histories from Theorem~\ref{ch09:step-theorem}. We then describe the
additional quantitative comparison needed for convergence of these
histories. Its abstract consequences are proved here. The distinction
between that comparison and ordinary continuity remains explicit;
in particular we do not assert analytic dependence on the tuned
endpoint. The elementary contraction arguments are applications of
Banach's principle~\cite{Banach1922}; we give the estimates specific to
the oppositely directed inverse coupling and activity recursions.

\section{What is fixed before choosing the scales}
\label{ch10:choices}

A stage $j\ge1$ specifies finitely many derivative orders, support
exponents, path widths and coefficient bounds. These are numerical
bounds for the representations of Chapter~\ref{ch04:blocking}, not a
new collection of physical theories. Stages are nested: the conditions
of stage $j+1$ include those already required at stage $j$. At a fixed
stage there are only finitely many permitted \emph{stage pairs} (the stages of the two scales of one step). Choose
constants for all these pairs simultaneously.

The order of choices is important. Fix the quadratic metric bounds,
path widths and classical bounds first. Choose the odd block ratio
large enough for the bound
\eqref{ch07:classical-budget}, the bound
\eqref{ch07:inherited}, and the margins in the support exponent for the local integral.
Only afterwards choose the lower bounds on the reference coupling and minimum auxiliary
periods. Enlarging such a lower bound does not change the block
kernel or the coefficient construction. The Wilson initialization can
be made smaller than any prescribed fixed fraction of its activity
bounds by this enlargement, as follows from the estimates in
Section~\ref{ch09:wilson}.

Let $C_j^{\rm drift}$ be a stage bound for the error in the downward inverse coupling step:
\begin{equation}\label{ch10:step-drift}
 \left|\beta_{m+1}-\beta_m-\kappa_G\log L_m\right|
       \le C^{\rm drift}_{J(m+1)}.
\end{equation}
Here $J(m+1)$ is the stage used for this step and
$\kappa_G=11C_{\gLie}/(24\pi^2)>0$. The normalization of the field-independent term is included in $C_j^{\rm drift}$.
Chapter~\ref{ch07:matching} fixes this coefficient with the same
normalization of the Lie algebra metric as the block integral.
Its error constant is fixed before the block ratio. Choose
nondecreasing odd integers $\mathscr L_j\ge3$ above all lower bounds on the block ratio at stage $j$,
and also require
\begin{equation}\label{ch10:drift-floor}
 \frac{C_j^{\rm drift}}{\log\mathscr L_j}
 \le \min\{\kappa_G/4,j^{-2}\}.
\end{equation}
Any further fixed logarithmic size bound $K_j$ can similarly be
required to satisfy $K_j/\log\mathscr L_j\le j^{-2}$.
Factors depending on the chosen $\mathscr L_j$ are covered by later lower bounds on the starting scales, not by redefining the constants used in
\eqref{ch10:drift-floor}.

Future applications may prescribe countably many finite numerical
requirements on the schedule, provided their analytic estimates have already been
established for a fixed stage. At stage $j$ take the maximum of the
first $j$ such requirements. This diagonal choice absorbs finitely many constants;
it does not establish an estimate merely because it has been listed.
In particular the comparison hypotheses below must be proved for the
same coefficient and activity construction before they can be used.


One particular later requirement can already be made explicit.
Fix the equation-adapted sections of
Theorem~\ref{ch02:equation-adapted} once and for all.
For $d\le j$, let $a_{dj}$ be their derivative constants at the
orders and spatial exponents of stage $j$, and let $q_d$ bound their
union stencil with $P_d^{\rm loc}$. These constants depend only on
fixed cover paths, finite polynomials and the stage orders;
Proposition~\ref{ch02:section-bounds} makes them independent of
the block ratio, period and cutoff. Fix the stage column bound
$C_j$, support exponent $\sigma_j$, and finite symmetry norm
$\mathcal O_d$ before choosing the ratio. With
\[
 Q_{dj}=2+e^{81\sigma_j}(a_{dj}+C_j+2q_d),
 \qquad \delta_d=(2\mathcal O_d^2)^{-1},
\]
include the additional lower bound on the block ratio
\begin{equation}\label{ch10:equation-floor}
 \mathscr L_j^{1/8}\ge
       \max_{d\le j}2(1+C_j)Q_{dj}\delta_d^{-1}.
\end{equation}
For the bounded coarse polynomial trials also include
$\mathscr L_j^{1/8}\ge8(1+C_j)Q^{\rm coarse}_{dj}$ for $d\le j$,
where $Q^{\rm coarse}_{dj}$ is their small-field chart bound, fixed before
the ratio. On that chart the indicator of the relaxed regularity condition is one. Its
complementary global value is controlled by the same later
square-root large-field bound condition described below.
Here $C_j$ is the prescribed bound, not a norm recomputed after
choosing $\mathscr L_j$. No radius at the ultraviolet cutoff enters this lower bound:
the raw stencil is fixed; the requirement at the ultraviolet cutoff that its cover
be at most $81$ is imposed later by increasing the lower bound on the
reference coupling and hence the available localization radius.
The remaining value norms at the ultraviolet cutoff contain inverse powers of the
small-field threshold. The lower bounds on the starting scales require their product with the
square root of the large-field bound at the cutoff to be at most one,
as well as the fixed-stencil support conditions and the conditions on the number of insertion points per coarse site.
They are absorbed by choosing sufficiently large lower bounds on the reference coupling,
not included in $Q_{dj}$. Those lower bounds on the starting scales also impose
$\widehat\beta_N\ge a_{d,J(N)}$ and all reserve inequalities at the cutoff
for $d\le J(N)$. Since only finitely many degrees are checked at a
stage, this is a noncircular addition to the schedule. Its use is
the near-identity comparison in
Section~\ref{ch11:equation-calibration}.

\section{A slowly increasing stage schedule}
\label{ch10:schedule}

We construct explicitly the \emph{schedule}: the stage and the block ratio at every scale, including the initial scales.
Let $b_j,c_j$ be nonnegative integers specifying respectively the
minimum starting scale and a requested length for a \emph{stage plateau}, a maximal interval of scales on which the stage is constant.
If $\jmath_j$ is the largest smoothing order used in stage $j$ or a
permitted stage pair, include
$b_j\log3\ge2\log(\jmath_j+7)$ for $j>1$. This ensures that
the full dependence radius of a probability-kernel pullback shrinks in
continuum units, including the variables that choose its local gauge transformations.
Starting with $A_1=0$, define
\[
 k_j=c_j+2A_j+1,\qquad
 A_{j+1}=A_j+b_{j+1}+2k_j+1,
 \qquad J(m)=j\quad(A_j\le m<A_{j+1}).
\]
Then $J$ is nondecreasing, tends to infinity, and increases by at
most one in a single step. If $J(m+1)=j>1$, then $m\ge b_j$.
Moreover the interval
$[k_j-\lfloor k_j/2\rfloor,k_j+\lfloor k_j/2\rfloor]$
lies in the $j$th plateau: its lower endpoint is at least $A_j$,
and its upper endpoint is less than $A_{j+1}$.

Set
\begin{equation}\label{ch10:scales}
 L_m=\mathscr L_{J(m+1)},\qquad
 \tau_0=0,\quad \tau_{m+1}=\tau_m+\log L_m,
 \qquad a_m=e^{-\tau_m}.
\end{equation}
Thus $a_m$ is the lattice spacing, and $a_m^{-1}$ is an integer.
Monotonicity gives
\[
 m\log3\le\tau_m\le m\log\mathscr L_{J(m)}.
\]
Choose a prescribed value $\bar\beta>0$ of the coupling at unit scale, the \emph{coarse endpoint}, and reference couplings
\begin{equation}\label{ch10:reference}
             \widehat\beta_m=\bar\beta+\kappa_G\tau_m.
\end{equation}
The lower bound of the first stage on the reference coupling is satisfied by $\bar\beta$. Every subsequent
such lower bound $d_j$ is satisfied by increasing its starting scale until
$\widehat\beta_m\ge d_j$. We may simultaneously require
\begin{equation}\label{ch10:late-entry}
 \log\mathscr L_j\le\tau_m/j,\qquad
 V_j(\tau_m)\le e^{\tau_m},
 \quad J(m+1)=j>1,
\end{equation}
Also require, for the original support radius and its per-site
normalization exponent, $B_{m+1}^2\le2B_m^2$ and
$(L_m B_m^2)^4q_{m+1}^{\rm site}/\tau_{m+1}\le1/j$ at stage $j$.
Here $q_{m+1}^{\rm site}$ is bounded by
$C_j(1+\log\widehat\beta_{m+1})$ using the Haar-ball normalization
argument. At a fixed stage the numerator is a fixed polynomial in
$\log(2+\tau_m)$, so this is a finite lower bound on the starting scale. It will keep
replication costs subpower at a nearby sparse stopping scale. For the
later grouping of finitely many local insertions, also require
$(\sigma_{m+1}+1)h_m/\tau_m\le1/j$, where
$h_m=16[1+L_m(B_m+B_{m+1}+1)^2]^4$ and $\sigma_{m+1}$ is the
prescribed support exponent. This is again a fixed polynomial in
$\log(2+\tau_m)$ at each stage. It ensures that moving a finite
number of insertion points contributes only a subpower factor.
In \eqref{ch10:late-entry}, $V_j$ is any prescribed fixed polynomial bound for the required
lattice side. The same argument works for a subexponential $V_j$.
Choose $b_j$ so these inequalities hold whenever
$\tau_m\ge b_j\log3$. They then hold at every step of stage $j$,
not just at its first occurrence.

For comparison estimates requiring an interval of scales $w_j(s)=o(s)$ around
one scale of stage $j$, increase $c_j$ so that
$w_j(s)\le\lfloor k/2\rfloor$ whenever
$k\ge c_j$ and $k\log3\le s\le k\log\mathscr L_j$.
At $k=k_j$ the whole prescribed interval stays on the same plateau.
This is useful when a fixed finite calculation must be isolated from
changes in derivative order or support exponent.

Only finitely many steps belong to the first plateau. Choose an integer
$M$ large enough that the lattice side $R_p/a_m$ meets the required
lower bound at every step of this plateau for all $p\ge0$, where
\begin{equation}\label{ch10:periods}
                 R_p=4(M+p+1),\qquad p\ge0.
\end{equation}
On later plateaux \eqref{ch10:late-entry} ensures the required lattice side. These periods also contain
arbitrarily late multiples of every positive integer: for a given
integer $c$, take $p=c(M+z+1)-(M+1)$ with $z$ arbitrarily large.
Finally
\[
 \frac{\log L_m}{\tau_m}\longrightarrow0.
\]
This follows from \eqref{ch10:late-entry} outside the first plateau.
No limit involving $\tau_0^{-1}$ is intended.

\section{Tuning a finite history to the coarse endpoint}
\label{ch10:shooting}

The coupling band at scale $m$ is
$[\widehat\beta_m/2,2\widehat\beta_m]$. Define inner barriers
\[
 \ell_m=\bar\beta+\tfrac12\kappa_G\tau_m,
 \qquad u_m=\bar\beta+2\kappa_G\tau_m.
\]
Both lie in the band. If a step starts at $\ell_{m+1}$, its output
is at most $\ell_m-\kappa_G\log L_m/4$ by
\eqref{ch10:step-drift}--\eqref{ch10:drift-floor}. Starting at
$u_{m+1}$ gives an output at least
$u_m+3\kappa_G\log L_m/4$.

We use only the intermediate value theorem. If a continuous function
$f$ on $[x,y]$ satisfies $f(x)\le d<u\le f(y)$, take its first
point $w$ with value $u$ and its last point $z\le w$ with value $d$.
Then $f([z,w])\subset[d,u]$. An excursion below $d$ would force a
later crossing of $d$; an excursion above $u$ an earlier crossing of
$u$. When $d=u$, choose a single point of that level set. In
particular no monotonicity of the output inverse coupling is required.

\begin{theorem}[Tuned finite histories]\label{ch10:finite-history}
Assume that every block step used below satisfies Assumption~\ref{ch09:forest-assumption}, which is not proved in this monograph. Fix the common numerical choices above and the compatible finite
quadratic and path constructions of Chapters~\ref{ch05:gaussian}
and \ref{ch07:matching}. For every cutoff $N$ and period $R_p$ there
is a finite history starting at a Wilson density and ending with
$\beta_{N,p,0}=\bar\beta$. It has the admissible activity bounds and
exact block identities at every scale $0\le m<N$. Its inverse couplings
stay between $\ell_m$ and $u_m$.
\end{theorem}

\begin{proof}
At scale $N$ use the Wilson curve with bare inverse coupling in
$[\ell_N,u_N]$. Apply the continuous step of
Theorem~\ref{ch09:step-theorem}. The barrier inequalities and the
interval argument give a compact subinterval on which the output
stays between $\ell_{N-1}$ and $u_{N-1}$ and attains these endpoints.
Restrict the entire previously constructed history to that interval.
Restriction compatibility preserves the indices, activity bounds and
exact identities. Repeat downwards. At scale zero the barriers
coincide at $\bar\beta$. For $N=0$ use the Wilson family directly.
Every input to an applied step lies in its band; no continuation
through a forbidden inverse coupling interval has been used.
\end{proof}

Choose one such history for each $N,p$, retaining its path arrays,
local integral choices (made for every forest and averaged over the forests as in Section~\ref{ch09:forest-average}) and output representations. The
choice need not be unique. The paths fixed before the period are
kept identical when the period changes. The exact identity is
\[
 \operatorname{dens}(\rho_{N,p,m})
 =\mathcal B_m^{\beta_{N,p,m+1}}
                  \operatorname{dens}(\rho_{N,p,m+1}).
\]
It is an identity of the densities, not merely of their representations, including the additive
normalization constant. That constant is not asserted to converge.

The same construction reaches any coarse endpoint
$|b-\bar\beta|\le h$, where
$h=\min\{\bar\beta/2,\kappa_G\log L_0/4\}$.
Descend to scale one and use the last barrier crossing to reach $b$.
This proves existence in a common endpoint interval, but neither
uniqueness of the bare coupling nor continuity of its choice.

\section{Uniform logarithmic running}
\label{ch10:running}

Telescoping from the fixed endpoint gives
\begin{equation}\label{ch10:running-error}
 |\beta_{N,p,m}-\widehat\beta_m|
 \le\sum_{i<m}C^{\rm drift}_{J(i+1)},\qquad N\ge m.
\end{equation}
For each $\epsilon>0$, choose $r$ so that
$J(i+1)^{-2}\le\epsilon/2$ for $i\ge r$. The right side is at most
\[
 \sum_{i<r}C^{\rm drift}_{J(i+1)}+\tfrac\epsilon2(\tau_m-\tau_r).
\]
The finite initial sum is eventually at most $\epsilon\tau_m/2$.
Consequently
\begin{equation}\label{ch10:running-limit}
 \beta_{N,p,m}=\bar\beta+\kappa_G\tau_m+o(\tau_m)
\end{equation}
uniformly in $N\ge m$ and $p$. In particular the bare Wilson
inverse coupling obeys $\beta_{N,p,N}/\log(a_N^{-1})\to\kappa_G$ uniformly
in $p$. This conclusion does not use convergence of histories at a
fixed scale. A bounded stage error is not summable by itself; the
small ratio in \eqref{ch10:drift-floor} is essential.

By using a smaller fixed constant than $\kappa_G/4$ in that lower bound,
the ratios $\beta_{N,p,m}/\widehat\beta_m$ can all be kept in an
arbitrarily prescribed fixed neighborhood of one. We henceforth
choose it small enough that
\[
 17/20\le\sqrt{\beta_{N,p,m}/\widehat\beta_m}\le23/20.
\]
For example an error at most $\kappa_G\log L_m/10$ at each step
implies a relative error at most $1/10$ after summation. This also
covers the initial indices, including $m=0$.

\section{The coefficient comparison before the activity comparison}
\label{ch10:path-stability}

A representation by transported curvature terms is not unique.
For comparisons we fix the linear curl division and the choice of coordinate contours
of Chapter~\ref{ch07:matching} once on the common
index sets of the coefficients. An alternative is the orthogonal
representative modulo the kernel of all periodic Taylor coefficients;
then the orthogonal complement must be taken inside the admissible
subspace, so that its invariance and moment constraints are preserved.
Equality of the action on one torus is insufficient to identify its
coefficient array.

At a fixed stage enlarge each incoming classical or loop coefficient bound
by a factor $64$, one family at a time, and choose the lower bound on the block ratio
for that enlarged bound. The bounds in
\eqref{ch07:classical-budget} and \eqref{ch07:inherited} then retain
unenlarged output bounds. This enlargement does not concern the
positive quadratic metric, which is kept fixed. The fixed-order
calculation at the critical configuration is valid after reducing its local analytic
neighborhood if necessary.

With quadratic data fixed, the outgoing cubic coefficient is affine
in the incoming cubic coefficient. With the cubic coefficient also
fixed, the same holds for the quartic coefficient. Finally the
quadratic loop coefficient is affine in the incoming loop. The
implicit-differentiation formula \eqref{ch07:uniform-stationary-jets} proves the
first two assertions: its exchange term is fixed in the quartic
comparison. The inherited Hessian and determinant formula prove the
third. Fixed linear coefficient functionals preserve these identities.

Normalize each coefficient norm by its stage bound. If an affine
map $T$ sends the admissible ball of radius $64$ into the output ball
of radius one, then for $\|H\|,\|K\|\le1$, writing
$\delta=\|H-K\|>0$, the admissible extrapolation
$K_*=K+32(H-K)/\delta$ has norm at most $33$ and gives
\begin{equation}\label{ch10:affine-contraction}
 \|T(H)-T(K)\|\le\delta/16.
\end{equation}
The extrapolation preserves all linear symmetry and moment conditions.
This is an affine estimate, not an assertion that the full nonlinear
step is affine.

Proposition~\ref{ch05:horizon-limit} gives convergence of the
quadratic kernels and free transports at every fixed scale, with a
strictly stronger exponential width than the required width. Choose
the direct coefficient strictly below its uniform positive metric
gap, as required there.
Their fixed-window inverses converge by the resolvent identity and the
common inverse bound. In each successive coefficient comparison the
change in these preceding data contributes an error
$e_{N,N',i}\to0$ at fixed $i$. The current coefficient discrepancy
therefore satisfies
\[
 \delta_i\le\tfrac14\delta_{i+1}+e_{N,N',i}.
\]
For $k>m$ and $N,N'>k$, iteration gives
\begin{equation}\label{ch10:path-truncation}
 \delta_m\le2(1/4)^{k-m}
             +\sum_{i=m}^{k-1}(1/4)^{i-m}e_{N,N',i}.
\end{equation}
First send $N,N'\to\infty$ at fixed $k$, then $k\to\infty$.
This proves the Cauchy property in succession for cubic, quartic and
loop coefficients, for the quadratic kernels of the nested blocking sequence and the fixed common coefficient functionals.

There is a useful precise version of the strict width margin. If coefficients
$c_N(\gamma)$ on a common index set converge coordinatewise and
$\sum_\gamma |c_N(\gamma)|e^{w'\ell(\gamma)}\le C$ with $w'>w$,
then their tails at width $w$ and length greater than $r$ are bounded
by $Ce^{-(w'-w)r}$. If only finitely many indices have
bounded length, finite-coordinate convergence and this tail bound
prove convergence at width $w$. For a Hilbert encoding instead, one
must verify summability of the encoding weights on that index set;
a bound on individual coefficients alone does not bound the number of
paths of a given length.

\section{The stronger local comparison required for cutoff limits}
\label{ch10:comparison-input}

Fix a period $R_p$ in continuum units. At scale $m$ let $\mathcal X_{m,p}$ be
the completion of the activity arrays in their normalized
weighted derivative and value norm. Zero seminorm arrays are
identified. Equality of their evaluated partition functions is not
this equivalence relation. Index sets, event indicators, supports and prescribed-site
rules are fixed at this scale independently of $N$. Families are extended
by zero to absent indices.

To state the comparison, set
\[
 \chi_m=\widehat\beta_m^{-1/32},\qquad
 \mathfrak b_m=\widehat\beta_m^{31/32},\qquad
 \lambda_m=\chi_m^{-1}\sqrt{\beta_m/\widehat\beta_m},\qquad
 b_m(\lambda)=\widehat\beta_m(\chi_m\lambda)^2.
\]
The symbol $b_m(\lambda)$ in this section is the inverse coupling coordinate
map; it is unrelated to the penalty parameter $b$. On the narrow band
$17/20\le\chi_m\lambda\le23/20$ we have
\begin{equation}\label{ch10:inverse-coordinate}
 \mathfrak b_m|\lambda-\lambda'|\le
               |b_m(\lambda)-b_m(\lambda')|.
\end{equation}
The extra factor in the right side is
$\chi_m(\lambda+\lambda')\ge17/10$.

Write the step maps at cutoff $N$ as
\[
 x_m=F_{m,N}(\lambda_{m+1},x_{m+1}),\qquad
 q_m=Q_{m,N}(\lambda_{m+1},x_{m+1}),
 \quad
 \beta_m-\beta_{m+1}-\alpha_{N,m}=q_m.
\]
Here $x_m$ is the activity array, the family of localized terms and large-field activities, and $q_m$ is the coefficient $\Delta$ of the coupling constant
renormalization of the step. It is not the vacuum energy renormalization. The coefficient
$\alpha_{N,m}$ fixed by the matching is the one already used in the exact step.

The following strengthens the continuous finite-step statement of
Theorem~\ref{ch09:step-theorem}. Its proof uses the additional numerical reserves already allowed
in the scale choices.
\begin{proposition}[Quantitative comparison on the common data]
\label{ch10:stability-assumption}
The step maps are defined for every prescribed narrow-band inverse coupling
sequence and every allowed activity array of norm at most $1/8$ (here an allowed array is in particular equivariant in the sense of Definition~\ref{ch09:equivariant}, a linear condition with fixed bijections of labels; the maps are those of Theorem~\ref{ch09:step-theorem} and are therefore subject to Assumption~\ref{ch09:forest-assumption}),
including zero terminal arrays. They preserve that ball, satisfy
$|Q_{m,N}|\le s_m/2$, and, on the same cutoff, obey
\begin{align}\label{ch10:lipschitz}
 \|F_{m,N}(z)-F_{m,N}(z')\|&\le\tfrac1{16}\|z-z'\|,\\
 |Q_{m,N}(z)-Q_{m,N}(z')|&\le4s_m\|z-z'\|,\notag
\end{align}
where $\|z-z'\|=\max\{|\lambda-\lambda'|,\|x-x'\|\}$.
The strengthened Wilson terminal arrays have norm at most $1/8$.
The numerical choices satisfy $\sum_{i<m}s_i\le \mathfrak b_m/128$.
The paths and $\alpha_{N,m}$ converge at each fixed scale with the
strict width margin above. Finally, for two sequences of admissible
inputs at different cutoffs, common limits of their scalar coordinates
and Cauchy activity arrays with difference tending to zero imply
output differences tending to zero, in both the activity norm and
the coupling constant renormalization. This last assertion applies also to every
zero-terminal-array history with prescribed inverse couplings.
\end{proposition}

The last condition compares changing maps $F_{m,N}$, not just changing
inputs to one fixed map. Uniform contractions by themselves do not
imply it. For example $F_N(x)=(-1)^N/10$ has Lipschitz constant zero
and no cutoff limit. Nor does smallness of the output imply the
small Lipschitz constant in \eqref{ch10:lipschitz}.

\begin{proof}[Proof of Proposition~\ref{ch10:stability-assumption}]
We derive the contraction from finite parameter evaluations (finite linear combinations of values along a real line), bound those
evaluations using the local expansion, sum the errors over scales, and
finally prove continuity as the cutoff is removed.

\begin{enumerate}[label=\emph{Step \arabic*.},ref=\arabic*,leftmargin=*,align=left]
\item \emph{Contraction from finite parameter evaluations.}
One sufficient quantitative estimate uses only real parameter lines.
For a line $z+tv$ of admissible inputs, $|t|\le1$, suppose every finite
choice of weights $a_i$ and points $t_i$ satisfies
\begin{align}\label{ch10:finite-tests}
 \left\|\sum_i a_i F(z+t_i v)\right\|
 &\le d\sup_{n\ge0}\left|\sum_i a_i(t_i/r)^n\right|,\\
 \left|\sum_i a_i Q(z+t_i v)\right|
 &\le f_m\sup_{n\ge0}\left|\sum_i a_i(t_i/(2r))^n\right|,
 \qquad r\ge2.\notag
\end{align}
These are estimates of the arrays in their derivative norm,
not just of the value of a partition function. They imply an
absolutely convergent one-variable expansion. Indeed the map
$((t/r)^n)_{n\ge0}\mapsto F(z+tv)$ extends linearly and boundedly
from its span in $c_0$ to its closure. That closure contains every
coordinate vector: subtract the first $n$ coordinates, divide by
$(t/r)^n$, and let $t\downarrow0$. Applying the bounded map to the
geometric expansion gives coefficients of norm at most $d$.
Consequently
\[
 \|F(z+tv)-F(z)\|\le d|t|/(r-|t|).
\]
The identical argument applies to $Q$. This short bounded-operator
argument records the precise information in the parameter evaluations; it is not an
inference of analyticity from continuity.

For two inputs at distance $0<h\le1/4$, use the line with direction
$(z'-z)/(4h)$ and evaluate at $t=4h$. Activity norm along this line
is at most $1/8+1/4<1$. An inverse coupling reserve from the narrow band to
$33/40\le\chi_m\lambda\le47/40$ makes the whole line admissible
when $\chi_m\le1/128$. The resulting bounds are at most $4dh$ and
$4f_mh$. If $h>1/4$, use the separate output bounds; if the distance
is zero, use the same estimate with a vanishing seminorm direction
and let the parameter tend to zero. Thus $d\le1/256$ and
$s_m=2f_m$ imply \eqref{ch10:lipschitz} with room to spare.
The finite-evaluation inequalities require the analytic estimates on the
local expansion established in the next step.

\item \emph{Parameter estimates from the local expansion.}
\label{ch10:parameter-supplier}

Here is the parameter estimate for the local factors. It uses
the protected change of variables and real Gaussian majorants of
Chapter~\ref{ch06:large-fields}, retaining their fixed derivative
order and margins in the support exponent. The incoming array is affine along the
parameter line. The path, event indicators, held fine reads and coefficient
projection are fixed. In particular the finite low-order collection in
Chapter~\ref{ch09:renormalization} is chosen on the compact inverse coupling
band independently of incoming activity arrays. The estimates below
therefore compare the same output construction on different lines.
For the estimates one uses the explicit local expansion on the larger
affine activity domain. Its polynomial products and convergent logarithmic
series define the array and local scalar outputs there even when an
extrapolated array has not been identified with a positive density.
Positivity is required for the histories of densities, not for this algebraic
extension of their local comparison formulas.

Put $y=\widehat\beta_{m+1}$, $R=y^{1/256}$ and
$\delta=y^{-1/128}$. The inverse coupling is a square of an affine function
of the line parameter. The relative change on $|t|\le6R$ is
$O(Ry^{-1/32})=O(y^{-7/256})=o(\delta)$. Choose a slightly smaller
reference inverse coupling $\beta_0$ so that
\[
 \beta(t)=\beta_0\ell(t)^2,\quad \ell(t)=a+ct,
 \qquad 1\le a-6R|c|\le a+6R|c|\le1+\delta.
\]
All these real inverse couplings remain in the reserved band. A compensating
rescaling keeps each fine field read fixed: the Gaussian coordinate is rescaled
by $\ell(t)$ while its fluctuation coefficient is divided by
$\ell(t)$. The indicators of large fluctuation-field events and measurable incoming values
are consequently not differentiated in this parameter.

For a parameter series use the norm
$|f|_R=\sum_{n\ge0}|f_n|R^n$. Its product estimate is the ordinary
Cauchy-product inequality. A real local energy factor, including its
coordinate prefactor, has the form
\[
 J(t,U)=A(U)\ell(t)^n\exp\{(1-\ell(t)^2)\Gamma(U)\}.
\]
The energy $\Gamma$ can have either sign. It cannot be replaced by
its absolute value before estimating the exponential.
For $q=\Gamma(U)$ the exponential series norm is at most $e^{H(q)}$,
where
\[
 H(q)=(1-a^2)q+2aR|cq|+R^2c^2|q|.
\]
Choose the sign of $\xi$ so that $q c\xi\le0$ and take
$3\le|\xi|\le4$. Since $a\ge1+6R|c|$, direct expansion gives
$H(q)\le(1-\ell(R\xi)^2)q$. After $j$ coarse-field derivatives,
factoring out $\ell(t)^n e^{(1-\ell(t)^2)q}$ leaves a polynomial
of degree at most $2j$. Interpolate that polynomial at $2j+1$
distinct such real nodes. Their scaled Lagrange polynomials have a
fixed coefficient norm. Thus real jet bounds $C d^i$ at those nodes
imply
\begin{equation}\label{ch10:real-node-bound}
 |D^iJ(\cdot,U)|_R\le K_j(1+\delta)^n C d^i,
                   \qquad i\le j.
\end{equation}
The node sign is chosen only for a pointwise bound and is never
differentiated. The same argument applied to the Gaussian density
ratio in $D$ local real coordinates yields
\[
 \left|\ell(t)^D e^{(1-\ell(t)^2)Q(z)}\right|_R
 \le(1+\delta)^D,\qquad Q(z)\ge0.
\]
Only coordinates read by the local term enter $D$.

The large-field corrections, which are small, require a different estimate. Their
parameter dependence is polynomial of degree at most two; real
values and derivatives of their exponential minus one have bound
$v d^i$ at $-1,0,1$. For $v\le1/2$, taking the real logarithm at
these three nodes preserves a bound $C_jv d^i$. Polynomial
interpolation contributes a factor $C_jR^2$, and exponentiating in the coefficient
algebra gives a bound $C_jR^2v d^i$ when $C_jR^2v\le1$.
Affine activities with a prescribed large-field set have the same polynomial interpolation
bound without taking a logarithm. These steps require no
analyticity of the incoming activity as a function of the gauge
field beyond the already imposed finite derivative order.

For the primitive bound $v=y^{-7/16+o(1)}$, this radius factor
changes its exponent to
$-7/16+2/256=-55/128<-2/5$.
The local factor estimate \eqref{ch10:real-node-bound} retains its negative
large-field energy: its additional $(1+\delta)^n$ is absorbed into
the fixed constant per local factor. A local read set in support $S$ has
at most $C_L(1+B)^4 C_B(S)$ coordinates. Since $B$ is polynomial
in $\log y$, its Gaussian prefactor satisfies, for any fixed
$\eta>0$ and sufficiently large $y$,
\[
 (1+\delta)^{C_L(1+B)^4C_B(S)}\le e^{\eta C_B(S)}.
\]
This is absorbed by a margin in the support exponent, independently of the
volume. The joint probability of large fluctuation-field events and Gaussian score moments
are those already used for the common majorant; Cauchy--Schwarz
retains their gain from these events without assuming that they are independent.

We can therefore run the connected expansions of
Chapter~\ref{ch08:chapter} in this coefficient norm. A derivative
of order at most $j$ of a product with $k$ factors contributes a factor at most
$k^j\le2^{jk}$, absorbed by enlarging each vertex bound by $2^j$.
Polynomial support moments are absorbed by another fixed margin in the
support exponent. The first connected expansion has bound $y^{-2/5+o(1)}$;
the logarithmic expansion and spatial derivative weights consequently
have bound $y^{-3/8}$ at radius $R/2$, after increasing the lower bound on the starting scale. The low-order projection uses finitely many auxiliary
volumes of side $O(B)$ and supports coverable by a fixed number of
$B$-balls. Its polynomial factors have exponent zero. Hence the
raw projected regular family has finite-evaluation bound
\begin{equation}\label{ch10:broad-test}
 \mathcal N\left(\sum_i a_i H(t_i)\right)
 \le y^{-1/4}\sup_{n\ge0}\left|\sum_i a_i(2t_i/R)^n\right|.
\end{equation}
Here $\mathcal N$ is the full unnormalized regular derivative bound.
The margin between $-3/8$ and $-1/4$ absorbs the projection factors.
A connected covering number merely growing subpolynomially would not suffice
inside an exponential weight; its uniform bound is used here.

For completeness we explain the sharpening which makes this useful.
At each real point, the stronger one-step remainder estimate gives
\[
 \mathcal N(H(t))\le 3\varepsilon C_{\rm st}x^{-11/10}
                         +O(x^{-23/20+o(1)}),\qquad
 x=\widehat\beta_m,
\]
where the fixed fraction $\varepsilon$ can be made as
small as desired before choosing the lower bound on the reference coupling. This is the
small remainder estimate of Chapter~\ref{ch09:renormalization}
applied to this same collection.
Let a finite parameter evaluation have moment norm $M_q$ at a fixed radius $q$.
Interpolation at $k+1$ fixed real nodes separates its moments through
degree $k$ from the remainder. The interpolating evaluation functional has total
variation at most $U_{k,q}M_q$; the remaining evaluation functional has moment norm
at radius $v\ge q$ at most
$(1+U_{k,q})(q/v)^{k+1}M_q$. Applying the pointwise bound to the
first part and \eqref{ch10:broad-test} to the second yields
\[
 M_q\{U_{k,q}\sup_t\mathcal N(H(t))
             +(1+U_{k,q})(2q/R)^{k+1}y^{-1/4}\}.
\]
Take $k=300$. The last power is $y^{-365/256}$, strictly smaller
than $y^{-23/20}$. First choose $\varepsilon$ to absorb the fixed
interpolation constant, then increase the lower bound on the reference coupling. The
result is an arbitrarily small multiple of $x^{-11/10}M_q$.
The scalar projection contributes the squared inverse field threshold as a factor;
thus its bound has the form
\[
 f_m=C\varepsilon x^{-11/10}/t(x)^2,
 \qquad t(x)=x^{-1/2}e^{(\log\log x)^2}.
\]
The coupling constant renormalization is linear on finite parameter evaluations, and its fixed path contribution
is $O(|\alpha|t(x)^3)=o(x^{-11/10})$.

Arrays with a prescribed large-field set retain one required large-field factor throughout
the exact regrouping. Work with their coefficientwise majorants, not separately
chosen value bounds. A raw hard bound $a_h\le C e^{-cP(y)^2}$,
where $P(y)=e^{(\log\log y)^2}$, is negligible compared with any
fixed fraction $\phi$ of $e^{-P(x)^{1/5}}$. The soft correction has
series bound $a_s\le y^{-1/8}$. Put $p=y^{1/16}$ and
$\zeta=-\log(\phi p)\ge0$. Then
$I=e^\zeta a_h+2a_s\le(8p)^{-1}+2a_s$ eventually, the support
ball count times $I$ tends to zero, and
\[
 e^{-\zeta}(e^I-1)\le\phi(1/4+4pa_s)\le\phi.
\]
This is precisely the estimate of the exact regrouping, with the hard factor
retained once. Radius reduction by a fixed factor converts the
finite-evaluation coefficient bounds into one summable majorant at every
support. Indicators of the eligibility conditions are parameter independent and
preserve it. The result proves both inequalities
\eqref{ch10:finite-tests}, with arbitrarily small fixed $d$, for the
protected local construction. The inverse coupling bound above gives
$s_m=2f_m\le C_jx^{-1/10}$ needed for cumulative control.

This argument uses the real-node bounds for the same protected local
sectors on the reserved inverse coupling band. Those bounds follow by applying
the real estimates of Chapter~\ref{ch06:large-fields} with their
constants fixed on that band; the rescaling keeps all fine reads and
event indicators fixed. It adds a parameter estimate to that construction,
not a differentiation of measurable cutoff indicators.

\item \emph{Summation over scales.}

The cumulative scalar bound is compatible with the schedule.
Suppose at a fixed stage its bound has the proved form
$s(x)\le C_j x^{-1/10}$, where $x$ is the reference coupling.
For the stage increment $a=\kappa_G\log\mathscr L_j>0$ and $x\ge a$,
\[
 (x+a)^{31/32}-x^{31/32}
 \ge\tfrac{31}{32}2^{-1/32}a x^{-1/32}.
\]
Since $x^{-1/10}=o(x^{-1/32})$, a lower bound on the starting scale ensures
$s(x)\le((x+a)^{31/32}-x^{31/32})/128$.
Summation gives the required bound by telescoping, including all
initial scales once $\bar\beta$ exceeds the lower bound of the first stage.
The decay estimate for $s(x)$ is required before this numerical
argument; an arbitrary stage constant would not suffice.

\item \emph{Continuity for a common input.}
\label{ch10:cutoff-supplier}

Use now the quadratic and path coefficient limits with the strict
width margin established in
Section~\ref{ch10:path-stability}. We verify the cutoff-comparison
part of Proposition~\ref{ch10:stability-assumption} on the
allowed inputs. All assertions in this paragraph are at one fixed
scale and period; constants need not be uniform in the scale.

For a finite collection of elementary indices and bounded path lengths,
the fine-coordinate dimension is finite. Write each integral in the
same physical coordinates at the two cutoffs. Convergence of the path
coefficients with a width margin gives convergence of its finitely
many required derivatives, uniformly on the smaller local chart.
The inverse identity
$H_N^{-1}-H^{-1}=H_N^{-1}(H-H_N)H^{-1}$, the common positive Hessian
bound and uniqueness of the critical configuration give convergence of means, precisions
and jets of the critical configuration. Differentiated Gaussian densities are their
ordinary polynomial scores. On a sufficiently late tail their
precisions have a common smaller positive lower bound, so a fixed
Gaussian exponential times a polynomial dominates them. The protected
change of variables leaves measurable fine reads fixed. Thus the
same input activity discrepancy is bounded in its prescribed norm;
no derivative of a measurable indicator is introduced.

Use the physical fine coordinates also for the event indicators. At a fixed
scale their threshold is determined by the fixed reference coupling,
and the list of large fluctuation-field events, the index set of supports and the clipping rule
are independent of the ultraviolet cutoff. Thus the event indicator is the
same measurable function in both integrals. Although its expression
in a centered Gaussian variable would move with the mean or precision,
no such parameter-dependent expression is compared here. The local change of
density includes its Jacobian and leaves that physical event indicator fixed.
The coarse small-field event indicators are likewise identical. This avoids any
uniform assertion about the measure of a parameter-dependent threshold surface.

To obtain the suprema in the output derivative bound, choose the
protected Gaussian and chart estimates through one order beyond the
outgoing order. This is one additional derivative order for the
integration kernels, and does not require one additional derivative
of an incoming activity. Indeed in each differentiated integral the
protected coordinate change holds its entire fine activity read
fixed. An additional coarse increment differentiates the smooth
density, chart and Jacobian scores, while the old measurable factors
remain unchanged. The bounds of Chapter~\ref{ch06:large-fields} apply
at every fixed prescribed score order; increase that order before
selecting the stage constants. Direct path terms and the finite
local projection are smooth and satisfy the same extra order by
their exponential coefficient bounds. Thus this additional order is available
also on a plateau with equal incoming and outgoing activity derivative
orders. On each common small-field chart it gives equicontinuity of the
finite collection of output derivatives. A finite net on the compact
coarse configuration domain and dominated convergence then give
uniform convergence there. The estimates extend to its closure from
the protected neighborhood; multiplication by the same small-field event indicator
preserves the resulting supremum bound.

Next remove the finite-index restriction. Long paths have uniformly
small tails by the strict exponential width margin. Large connected
supports and high cluster multiplicities have uniformly small tails
by the margins in the support exponent and the attachment reserves of the polymer estimates.
Remaining multiplicities inherited from the input arrays are treated
by the Cauchy tail argument at the end of this chapter: a fixed late
array has a summable value tail, and the activity-norm discrepancy
bounds the tail of every later array. In a protected derivative the
input factors are kept at fixed fine reads, so these are value tails;
the differentiations fall on the controlled Gaussian and chart
factors. If an input derivative is explicitly present, its difference
is already controlled by the corresponding component of the input
norm. First truncate with these uniform errors, then compare the
finite integrals. This proves convergence in the output norm.

The coupling constant renormalization is a finite linear functional of the same
converging local derivatives and therefore also converges. Local
subtraction preserves this conclusion by linearity; the exact regrouping preserves
it by its convergent product expansion with a strict margin in the
site-norm bound. The finite auxiliary projection uses the same fields and
periods at both cutoffs. Thus this comparison concerns the
returned arrays and $Q_{m,N}$, not some other representation of their
densities.

The same proof applies when the terminal activity is zero and the
inverse couplings below it are prescribed in advance. At the initial auxiliary
scale the density is the positive bare path exponential, and all
activity bounds are zero. After each step its activity array has the
small bound proved above. Replacing the returned inverse coupling by the next
prescribed inverse coupling multiplies the represented density by a strictly
positive path exponential; its activity gas is unchanged. Therefore
positivity is retained, and the next step is an allowed input to
Theorem~\ref{ch09:step-theorem}. All conditions on the critical configuration, supports and cutoffs
depend on the common path and reference band, which have
not changed. In particular this induction supplies precisely the
histories with prescribed inverse couplings used in the convergence proof; it does not assume
that every abstract element of a completed norm ball represents a
positive density.

The quadratic limit of Proposition~\ref{ch05:horizon-limit}, the affine
coefficient comparison and the preceding local-integral argument give
the common-input continuity required here, including all zero-array
histories with prescribed inverse couplings. Together with the finite-evaluation
estimates and cumulative bound, this proves
Proposition~\ref{ch10:stability-assumption}. \qedhere
\end{enumerate}
\end{proof}

\section{Convergence of the coupled histories}
\label{ch10:fixed-scale}

Recall the schedule from Section~\ref{ch10:schedule}: $J(m)$ is the stage at scale $m$, $L_m=\mathscr L_{J(m+1)}$ is the block ratio, and $\tau_m=\log a_m^{-1}$. The reference coupling $\widehat\beta_m=\bar\beta+\kappa_G\tau_m$ is distinct from the tuned coupling $\beta_m$.

\begin{theorem}[Fixed-scale convergence]
\label{ch10:convergence}
Use the common trajectory above with its strengthened comparison
reserves and fixed coefficient functionals. If the coarse endpoints
converge to $b$, then for every fixed
$m,p$ their inverse couplings and activity arrays converge along the entire
cutoff sequence. Their coefficients $q_m$ of the coupling constant renormalization converge as well.
The limits depend on $b$ and the fixed common trajectory, and not
on the choices made in tuning the bare coupling. No convergence of additive vacuum
normalizations is asserted.
\end{theorem}

\begin{proof}
First prescribe two inverse coupling sequences having the same limit at each
fixed scale. At a scale $k>m$ replace the activity array by zero and
run the same maps downwards using the prescribed inverse coupling sequence, without
imposing the scalar recursion on this auxiliary history. The
hypothesis on its domain is needed here. Repeated use of
\eqref{ch10:lipschitz} gives
\begin{equation}\label{ch10:activity-truncation}
 \|x_m-x_m^{[k]}\|\le(1/16)^{k-m}.
\end{equation}
Figure~\ref{ch10:fig-cutoff-histories} shows the common interval
of scales on which the two truncated arrays are compared.

\begin{figure}[!hbp]
\centering
\begin{tikzpicture}[font=\small, x=1cm, y=1cm, >=Stealth,
    history/.style={ymInk, line width=.9pt},
    truncated/.style={ymAccent, dashed, line width=.9pt}]
  \node[anchor=east] at (-.3,.66) {level};
  \foreach \x/\level in {0/0,2.8/m,6.2/k,8.8/N,10.4/{N'}}
    \node at (\x,.66) {$\level$};
  \foreach \x in {2.8,6.2}
    \draw[ymGrid] (\x,.42) -- (\x,-1.65);

  \foreach \y/\cutoff/\name in {0/8.8/N,-1.17/10.4/{N'}} {
    \begin{scope}[shift={(0,\y)}]
      \node[anchor=east] at (-.3,0) {cutoff $\name$};
      \draw[history,->] (\cutoff,0) -- (6.29,0);
      \draw[history,->] (6.2,0) -- (2.89,0);
      \draw[history,->] (2.8,0) -- (.09,0);
      \foreach \x in {0,2.8,6.2,\cutoff}
        \fill[ymInk] (\x,0) circle[radius=1.7pt];
      \draw[truncated,->] (6.2,-.35) -- (2.8,-.35);
      \draw[ymAccent,fill=white,line width=.8pt]
        (6.14,-.41) rectangle (6.26,-.29);
      \node[anchor=west,inner xsep=3pt] at (6.27,-.35) {$x_k^{[k]}=0$};
    \end{scope}
  }
\end{tikzpicture}
\caption[Comparing two cutoff histories at a fixed level]{Activity comparison at a fixed period $R_p$. The two solid
rows reach the same tuned inverse coupling $\bar\beta$ at level zero.
In each dashed run the activity is reset to zero at $k$, retaining
that row's inverse couplings and block maps down to $m$.
At $m$ the error from this replacement is at most $16^{-(k-m)}$ in
each row, uniformly in its cutoff.}
\label{ch10:fig-cutoff-histories}
\end{figure}


For fixed $k$, the truncated arrays converge by finite downward
induction and the common-input cutoff comparison. This statement
is Cauchy convergence: apply that comparison to arbitrary pairs of
subsequences. Then choose $k$ so the two errors in
\eqref{ch10:activity-truncation} are small. Completeness gives a
common activity limit for both prescribed inverse coupling sequences, and common limits of the coupling constant
renormalizations follow from the same local comparison.

At each fixed scale the inverse coupling coordinate lies in a compact real
interval. A diagonal subsequence therefore makes all inverse coupling
coordinates converge; the preceding paragraph supplies activity
limits without compactness of their norm balls. Compare two such
limiting histories. The finite scalar recursion, summed from their
common coarse endpoint, gives
\begin{equation}\label{ch10:prefix}
 b_m(\lambda_m)-b_m(\lambda'_m)
                 =-\sum_{i<m}(q_i-q'_i).
\end{equation}
The limits of the common coefficients $\alpha_{N,i}$ fixed by the matching cancel.
There is no infinite sum or ultraviolet boundary term.

The estimates \eqref{ch10:lipschitz} pass to the two limits. To
justify this without extending a map arbitrarily to the completed
space, at each finite cutoff run the second limiting inverse coupling sequence
with a zero terminal array. The first paragraph identifies its limit
with the second history. Compare it with the first history at the
same cutoff and pass to the limit.
Put
$d_m=\max\{|\lambda_m-\lambda'_m|,\|x_m-x'_m\|\}$.
The activity discrepancy is at most $1/4$. By
\eqref{ch10:inverse-coordinate}, \eqref{ch10:prefix}, and the
individual scalar bounds, its inverse coupling discrepancy is at most
$\mathfrak b_m^{-1}\sum_{i<m}s_i\le1/128$. Hence
$D=\sup_m d_m\le1/4$. The sharper scalar comparison now gives
\[
 |\lambda_m-\lambda'_m|
 \le4\mathfrak b_m^{-1}\sum_{i<m}s_iD\le D/32,
 \qquad \|x_m-x'_m\|\le D/16.
\]
Thus $D\le D/16$ and $D=0$. Every subsequence has a further
subsequence with these same limits. A subsequence staying a positive
distance away at one fixed scale would contradict this fact, so the
whole sequence converges. Finally $b_m$ is continuous, giving the
asserted convergence of the physical inverse couplings.
\end{proof}

\section{Identifying the analytic comparison on the selected family}
\label{ch10:identification}

For clarity we record what the additional input requires of the local
integral, so it can be checked independently of the preceding abstract
argument. The finite path comparisons must use the same fixed
coefficient functionals and rooted estimates as the history. Replacing the
Wilson normalization at the ultraviolet cutoff must leave every fixed lower-scale
construction unchanged eventually, or have an explicitly vanishing
comparison error. A new choice of equally bounded coefficients at
each occurrence would not suffice.

For a finite set of activity indices, write both local integrals in
common physical coordinates. Converging means and positive precisions,
with their uniform inverse and derivative bounds, give convergence
against a common Gaussian majorant. The protected change of variables keeps the physical fine event indicators fixed,
including their reference thresholds, and retains its common Jacobian. The common majorant of
Chapter~\ref{ch06:large-fields} concerns the inhomogeneous term and the event indicator;
it cannot be replaced by a bound for an unrelated flat Gaussian.
The coupling constant renormalization also needs convergence in its local
derivative norm. Convergence of the gas alone does not identify it.

The infinite index sets require a separate tail argument. For an
input array $f$, let $b_f(i)$ be the supremum of the absolute value of the
activity, including its event indicator. On a fixed finite torus the activity bounds
give $\sum_i b_f(i)\le C_m\|f\|$. A Cauchy sequence has uniformly
small value tails: compare with one late fixed member, truncate that
member's summable tail, and use
$b_f(i)\le b_{f-g}(i)+b_g(i)$. This does not assert density of finite
arrays in every derivative norm. Output derivative tails require the
larger margin in the support exponent in the connected expansion. The local
projection is finite and the sum of the exact regrouping has a strict convergence
margin, so these estimates, when proved on the common inputs, survive
both operations.

The auxiliary zero-array histories used in the proof require the same
checks. Their incoming activity bounds vanish, but their path,
critical configuration, cutoff geometry and differentiated Gaussian factors
do not vanish. One must use the same linearization of the critical configuration, clipping data and
critical-configuration map. Uniform implicit-function uniqueness identifies an
auxiliary map of this kind built from the same repeated steps with the selected one; merely invoking the
existence of some admissible critical-configuration map is insufficient.

The schedule,
finite-history tuning and logarithmic running follow from the finite-step theorem. The fixed-scale conclusion is
Theorem~\ref{ch10:convergence}. Steps~\ref{ch10:parameter-supplier} and
\ref{ch10:cutoff-supplier} in the proof of Proposition~\ref{ch10:stability-assumption}
establish the stronger local comparisons
from the protected expansion and converging background coefficients.
Proposition~\ref{ch05:horizon-limit} supplies the fixed-scale quadratic
limit for the nested blocking sequence. This completes the fixed-scale
construction, which remains distinct from subsequent convergence of
all continuum correlations.
